
\clearpage
\thispagestyle{fancy}
\noindent\begin{minipage}[b]{0.52\linewidth}\includegraphics[height=2.2cm]{logo-javeriana.png}\end{minipage}\hfill\begin{minipage}[b]{0.43\linewidth}\raggedleft\fontsize{14}{17}\selectfont Ingeniería de Sistemas\end{minipage}
\vspace{5mm}
\noindent\textbf{Historial de Cambios}\par\vspace{2mm}
\begingroup\small\setlength{\tabcolsep}{5pt}\renewcommand{\arraystretch}{1.2}
\noindent\begin{tabular}{|p{0.11\linewidth}|p{0.22\linewidth}|p{\dimexpr0.67\linewidth-6\tabcolsep-4\arrayrulewidth\relax}|}\hline
Fecha & Autor & Descripción\\[2mm]
25-Oct-2026 & Daniel Felipe Ramirez Vargas & Primera versión. Incorpora el contexto vigente del proyecto, las especificaciones de la plantilla, los aportes pertinentes que hemos realizado en trabajos pasados como el BOI, VPP y PAFP, y las indicaciones de la profesora y del director. Queda registrado en cada sección que aspectos siguen pendientes por resolver y se utiliza citación IEEE.\\\hline
27/10/2026 & Daniel Felipe Ramirez Vargas & Incorporación de glosarios, actualización de ficha de estudiantes y unificación de terminología FDS/WFDS.\\\hline
1-Nov-2026 & Andres Felipe Pinzon Marquez & Actualización de la Sección 3.2 (Construcción) para detallar la ejecución en el clúster HPC-Zine, la integración meteorológica, el modelo de costos y los criterios cuantitativos de contraste y cobertura. \\\hline
2-Nov-2026 & Andres Felipe Pinzon Marquez & Reestructuración formal de la Sección 3.2 (Construcción) dividida en Método, Actividades y Resultados Esperados; incorpora la adaptación de DAD, las 8 actividades secuenciales, la distinción de entregables formales (E2-E7) y la documentación interna de calidad bajo normas ISO/IEEE. \\\hline
2-Nov-2026 & Miguel Ángel Sarmiento Cuervo & Desarrollo y actualización de las secciones 4, 5 y 6. Revisión de compromisos institucionales y derechos patrimoniales; ampliación del marco teórico con fundamentos de simulación e infraestructura HPC-ZINE; investigación de proyectos similares nacionales e internacionales, comparación de alternativas y actualización y verificación de referencias bibliográficas.\\\hline
  2-Nov-2026 & Daniel Felipe Ramírez Vargas & Actualización del problema, solución y figura de agregación; objetivos, entregables y estándares; fuente de amenaza de incendios y metodología de Construcción y Semi-transición.\\\hline
2-Nov-2026 & Guillermo Andrés Aponte Cárdenas & Concretado de el análisis de impacto. Descripción de la fase metodológica 2: Inicio y de todos sus componentes; Método, actividades y resultados esperados.\\\hline
  
  \textcolor{black}{21-Nov-2026} & Daniel Felipe Ramírez Vargas & \textcolor{black}{• Secciones 1.3, 3 y 3.1.3. Unificación de la denominación de la fase adaptada como Semi-transición.\newline • Secciones 3.1.1 y 3.1.2. Corrección de las citas del PAFP y de las normas ISO/IEC/IEEE 29148 y 12207.}\newline \textcolor{black}{• Secciones 1.1.1, 1.1.2, 1.1.3, 1.1.4 y 2. Ajustes de redacción y conectores.}\newline \textcolor{black}{• Secciones 1.2.2, 1.3, 3.1.3 y 3.3.2. Unificación del nombre de la tabla de referencia cruzada.}\newline \textcolor{black}{• Secciones 3.1.2 y 3.2.2. Actividades por fase.\newline • Secciones 3.1.3 y 3.2.3. Resultados y relación con los objetivos específicos.}\\\hline
\end{tabular}\endgroup
\clearpage
